\section{什么是涌现，以及它不是什么}

\begin{definition}[涌现]
\textbf{涌现}是指这样一种高层有组织模式的出现：它已无法被有效描述为若干孤立局部事件的简单清单，
因为这一模式会反馈性地影响后续局部行为、改变未来成本，或成为一个本身具有解释价值的层级。
\end{definition}

\begin{definition}[框架形成]
\textbf{框架形成}是这样一个过程：重复输入、条件作用、环境耦合与历史性复用，共同生成一个相对稳定的行为或认知机制。
\end{definition}

\begin{remark}
本章使用“涌现”一词，是在一种有纪律的解释意义上，而不是把它当作任何复杂事物的魔法标签。
一种模式并不会因为它令人惊讶、难以处理或情感上重要，就自动成为涌现现象。
\end{remark}

\subsection{涌现不意味着什么}

对本手稿而言，涌现\emph{并不}意味着：
\begin{enumerate}[nosep]
  \item \textbf{神秘膨胀：} 仅仅因为微观细节不便追踪，就把某个过程称为涌现；
  \item \textbf{因果含混：} 仿佛一旦宏观模式出现，局部机制就不再重要；
  \item \textbf{无限制类比：} 仅仅因为某种生物或物理系统中存在涌现，就假定同一机制必然一一对应地支配某个具体人类常规；
  \item \textbf{词汇即证明：} 把“盆地”“相位”或“条件作用”之类系统语言，直接升级为关于日常行动性的证据，而缺少经验桥梁。
\end{enumerate}

\begin{discussion}
这一负面定义在策略上十分关键。
如果没有它，本章就会招致本应避免的批评：把“涌现”当成一种声望标签，而不是受约束的解释工具。
本书的主张更为狭义：重复的局部互动可以生成对历史敏感的组织；这种组织会进一步塑造什么更容易、什么更显著、什么更可能发生；因此，这个高层机制值得被命名。
\end{discussion}